\section*{Results}

\subsection*{Simulation validation: ArchaicPainter performance and ablation analysis}

We evaluated \AP{} across 50 independent simulation replicates
(msprime v1.3.3~\citep{baumdicker2022efficient}, 5~Mb per replicate,
20~CEU query haplotypes, 2\% Neanderthal admixture, seeds 42--91).
\AP{} with the standard bidirectional emission achieved a mean segment-level
F1 of $0.480 \pm 0.095$ (95\% CI $[0.454, 0.506]$; AUPRC $= 0.577$;
Figure~2).
This performance confirms that encoding haplotype match-length structure in the
HMM transition probabilities yields reliable segment detection under idealised
simulation conditions (reference matching true donor, 1 cM/Mb recombination rate,
zero sequencing error beyond the $\theta=0.01$ emission parameter).
\AP{} AUPRC of 0.577 indicates that posterior log-scores provide useful
segment rankings across operating points (though scores from the positive-only
model used on real data are approximate heuristics rather than calibrated
probabilities; see Methods).
Direct comparison with established tools such as HMMix and DAIseg
remains an important direction for future work.

\subsection*{Ablation analysis: reference fidelity is critical; segment merging is not}

Two ablation variants isolated the contribution of model components.

Disabling segment merging yielded F1 = $0.488 \pm 0.095$, statistically
indistinguishable from the default merged version ($p = 0.999$, Wilcoxon).
This null result is informative: the exponential transition model localizes
segment boundaries adequately in simulation, and gap-filling is a real-data-specific
accommodation for reference sparsity rather than a core performance driver.

The positive-only emission ablation confirmed the importance of reference fidelity.
Down-weighting mismatches reduced simulation performance to F1 $= 0.154 \pm 0.079$
versus $0.480$ for the bidirectional model ($p = 8.9 \times 10^{-16}$, Wilcoxon signed-rank test,
$n = 50$).
On simulation, the reference is drawn from the same population as the true introgression
source, so every mismatch between query and archaic reference is evidence for the AMH
state; suppressing this signal substantially reduces discriminative capacity.
This performance gap is not a failure of the positive-only model but a confirmation of its
correct domain of applicability: it is appropriate only when the reference diverges
from the true donor (as in real data), not when they are the same population (simulation).
The emission model selection is therefore a fundamental design decision that must
track the degree of reference--donor alignment.

\subsection*{\AP{} detects biologically coherent Neanderthal ancestry at per-haplotype resolution}

Applied to chromosome~21 (chr21:10--46.7~Mb, GRCh38; $\approx$36.7~Mb euchromatic region) from the
1000~Genomes Project Phase~3 GRCh38 release, \AP{} detected 187 introgressed segments
across 40~CEU haplotypes (mean ancestry fraction $1.10\% \pm 0.60\%$;
median segment length 56.7~kb) and 208 segments across 40~CHB haplotypes
($0.98\% \pm 0.52\%$; median 27.5~kb) (Table~2; Figure~3).
These fractions are broadly consistent with published estimates of $\sim$1--2\%
Neanderthal ancestry in non-African populations~\citep{prufer2014complete, sankararaman2014genomic}.
Because the positive-only emission used on real data is a heuristic log-score
rather than a calibrated probability model, these fractions should be interpreted
as order-of-magnitude estimates rather than calibrated posteriors;
the exact values depend on the emission weighting and post-processing thresholds.

Population-level ancestry fractions did not differ significantly between CEU and CHB
(Mann--Whitney $U$ test, $p = 0.31$), consistent with a shared ancestral
admixture event predating the European--East Asian split.
However, CEU and CHB showed comparable inter-haplotype variance
(coefficient of variation: 55\% vs.~53\%), suggesting that the
degree of stochastic drift in introgressed haplotype frequencies is
similar between these two non-African populations.
The established chr21:$\sim$14~Mb Neanderthal introgression hotspot was
detected in both populations~\citep{vernot2014resurrecting, sankararaman2014genomic, vernot2016excavating},
serving as a positive landmark for real-data evaluation.

\subsection*{\AP{} and IBDmix reveal complementary detection profiles}

Comparing \AP{} and IBDmix on chr21:10--20~Mb (10~Mb, 20~CEU individuals, LOD $\geq 3$)
highlighted fundamentally different operating characteristics.
IBDmix called 345 segments across 20 diploid individuals (17.3 per individual;
mean length 15.2~kb; median 11.5~kb).
The IBDmix \emph{diploid} ancestry fraction (total introgressed bp across both
haplotypes of all individuals, divided by $2 \times N_\text{ind} \times
\text{window\_bp}$) was 2.62\%.
\AP{} called 50 segments across 40 haplotypes (1.25 per haplotype; mean length
69.4~kb); the \emph{per-haplotype} ancestry fraction
(introgressed bp per haplotype $\div$ window bp) was 0.87\%.
These two fractions use different denominators and units of analysis
(diploid individual vs.\ phased haplotype) and should not be directly compared;
a rough diploid-equivalent for \AP{} can be obtained as $2 \times 0.87\% = 1.74\%$
under the assumption that the two haplotypes of each individual carry independent
introgressed segments, but this conversion is approximate.
Segment counts differed $\sim$7-fold and mean segment lengths $\sim$4.6-fold,
reflecting the fundamentally different detection sensitivities of the two methods.
In total, 31.0\% of IBDmix segments (107/345) overlapped at least one \AP{} segment;
conversely, 95.9\% of \AP{} segments in this region (47/49) overlapped at least one IBDmix segment.
This complementarity is architecturally principled: IBDmix uses site-level
LOD scores to achieve maximum recall of short tracts, while \AP{} requires
an extended haplotype match across many kilobases, trading short-segment
sensitivity for long-haplotype specificity and per-haplotype resolution.
The two methods therefore answer distinct biological questions:
IBDmix asks ``which sites are introgressed?''; \AP{} asks ``which haplotypes
carry long introgressed blocks and with what haplotype-level confidence?''\footnote{Throughout, \emph{confidence} refers to the Viterbi log-score under the positive-only
emission, not a calibrated posterior probability; see Methods.}

\subsection*{Three-state HMM achieves proof-of-concept Neanderthal--Denisovan source attribution}

Using a two-source simulation (20 replicates, seeds 100--119; 5~Mb;
2\% NEA + 1\% DEN admixture; matched simulated references),
the three-state \AP{} HMM detected both archaic ancestries above chance.
Neanderthal detection yielded F1 = $0.260 \pm 0.109$ with source attribution
accuracy $26.1\% \pm 13.5\%$; Denisovan detection yielded F1 = $0.188 \pm 0.105$
with attribution accuracy $15.2\% \pm 10.1\%$.
Cross-contamination was 11.9\% (NEA predicted overlapping true DEN)
and 19.3\% (DEN overlapping true NEA).
As a reference level, a null model that predicts archaic state randomly
in proportion to prior probabilities ($\pi_\mathrm{NEA} = 0.018$,
$\pi_\mathrm{DEN} = 0.002$) would place $\sim$99.9\% of calls in the NEA
state and thus achieve near-zero DEN recall; a uniform-random source
assignment conditional on a call being archaic would give $\sim$50\%
cross-contamination under the simulation's 2:1 NEA:DEN ratio.
The observed cross-contamination rates (11.9\% and 19.3\%) thus indicate
genuine discriminative capacity beyond chance.

The moderate attribution accuracies reflect two compounding biological constraints.
First, the $\sim$700-kya Neanderthal--Denisovan
divergence~\citep{prufer2014complete} limits the number of lineage-specific
alleles available for disambiguation.
Second, the low DEN admixture fraction (1\%) means some replicates contain
fewer than ten true Denisovan tracts, amplifying stochastic variance in per-replicate metrics.
These results are reported as a proof-of-concept of the three-state architecture;
the key contribution is demonstrating that a single HMM can simultaneously
attribute ancestry to two distinct archaic sources without post-processing,
providing a tractable path toward real-data Neanderthal--Denisovan disentanglement
in populations with high Denisovan ancestry (e.g., Papuans).

\subsection*{Near-linear computational scaling enables population-scale deployment}

\AP{} scales near-linearly with query panel size (Figure~4).
Per-haplotype runtime stabilized from $0.051$~s at $n = 10$ to $0.114$~s
at $n = 1{,}000$ haplotypes (2.2$\times$ increase over a 100$\times$ panel-size increase)
on a single CPU core.
The empirical scaling exponent was 1.21, close to the theoretical $O(n)$.
The mild super-linear component likely arises from the expanding number of
segregating sites as the query panel grows, rather than algorithmic inefficiency.
Extrapolating to a full genome ($\approx$3 Gb) for $n = 1{,}000$ haplotypes:
using bp-based scaling from the simulation benchmark yields $\sim$19 cpu-hours;
using per-site scaling at the chr21 informative-site density (27 SNPs/Mb, typical
for a single archaic reference) yields $\sim$1 cpu-hour, since Vindija's callable
sites are far sparser than simulation segregating sites.
Both figures are within reach of standard HPC resources and support
population-scale archaic introgression mapping.

\subsection*{Introgressed chr21 regions overlap adaptive immunity and neurological disease genes}

Intersecting the 52 merged introgressed regions (9.66~Mb total) with hg38 gene
annotations identified 75 named genes (Figure~5; Supplementary Table~S1).
Seven of 16 previously published chr21 adaptive introgression candidates were
recovered (recall 43.8\%), including \emph{DYRK1A} (neurodevelopment and autism;
CEU-specific), \emph{RUNX1} (hematopoiesis and leukemia risk; CHB-specific),
\emph{B3GALT5} (blood group antigen glycosylation; CHB-specific), and
\emph{LIPI} (triglyceride hydrolysis; shared between both populations).
We note that the merged introgressed regions span approximately 9.66~Mb of the
36.7-Mb analysis window ($\approx$26\%), so some gene overlap is expected
by chance; formal statistical enrichment (permutation tests against matched
genomic regions controlling for gene density and mappability) would be required
to establish significance beyond background expectation.
These results should therefore be interpreted as an illustrative demonstration
that \AP{}'s haplotype match-length signal recovers previously reported
candidate loci in an unsupervised manner, rather than as a rigorous
test of adaptive introgression.

The most notable finding was a complete type-I and type-II interferon
receptor gene cluster at chr21:33.2--33.5~Mb: \emph{IFNAR1}, \emph{IFNAR2},
\emph{IFNGR2}, and \emph{IL10RB}, all falling within \AP{}-predicted
introgressed regions predominantly in CEU haplotypes (Figure~5).
This cluster encodes the receptor subunits for interferons $\alpha/\beta$,
$\gamma$, and interleukin-10, defining the molecular entry point for innate
antiviral and anti-inflammatory signaling.
The detection of this cluster by \AP{} from haplotype match-length patterns alone
is consistent with published evidence that archaic introgression at this locus
contributes to elevated interferon response capacity in European
populations~\citep{dannemann2017contribution}, and illustrates that the
method's match-length signal can rediscover known functionally annotated
archaic introgression candidates in an unsupervised manner.

Population-specific gene overlaps further highlight divergent adaptive landscapes:
CEU uniquely spanned \emph{BACE2} (Alzheimer-pathway beta-secretase, chr21 paralog
of the BACE1 drug target), \emph{UBASH3A} (T-cell negative regulator with
genome-wide significant GWAS associations for type~1 diabetes and rheumatoid arthritis),
and the neuronal wiring gene \emph{DSCAM} (Down Syndrome Cell Adhesion Molecule);
CHB uniquely included \emph{ADAMTS5}, a metalloproteinase with replicated GWAS
signals for osteoarthritis, and \emph{RCAN1} (regulator of calcineurin-1,
associated with cardiac hypertrophy and Down syndrome cognitive phenotypes).
These population differences, combined with the known geographic structure of
adaptive introgression~\citep{vernot2016excavating}, suggest candidate loci
for differential archaic haplotype retention in European versus East Asian
populations; formal tests of positive selection at these loci (e.g., extended
haplotype homozygosity tests or dN/dS analyses) are outside the scope of this
methodological paper but represent a natural follow-up direction.
